% !TEX program = xelatex
\documentclass{article}

\usepackage{mathtools}
\usepackage{amssymb}
\usepackage{unicode-math}
\setmonofont{DejaVu Sans Mono} % monospaced font with extensive unicode support for verbatim
\usepackage{stmaryrd}
\usepackage{booktabs}
\usepackage{tikz}
\usepackage{forest}
\usepackage{logicproof}
\usepackage{ebproof}
\usepackage{fitch}
\usepackage[capitalize,noabbrev]{cleveref}

%%%%%%% Some macros to help with the typesetting %%%%%%%%%%
\usepackage{logic-macros}
%%%%%% End of local macros %%%%%%%

\title{Logic Homework 2}
\author{Jenny Vermeltfoort\\
  s3787494\\
  Informatica%Informatica, I&E, Bioinformatica, Wiskunde, ...
}

\begin{document}

%% Use var symbols by default
\renewcommand*{\phi}{\varphi}
\renewcommand*{\epsilon}{\varepsilon}
\DeclareMathAlphabet\mathzapf       {T1}{pzc} {mb} {it}

\maketitle

\section*{Question 1}

%[[\phi]]v(x) = [[\psi]]v(y,z) = 1 iff [[Ay.P(y)]]v(y) = 1 or [[Ez.R(z)]]v(z) = 1
%                                  if min{[[P(y)]]v[y->a] | a \in |\mathzapf{M}|} = 0, then by our assumptions  [[R(x)]]v[x->a] = 1 for some a \in |\mathzapf{M}|
%                                  iff min{[[P(y)]]v[y->a] | a \in |\mathzapf{M}|} = 1, then [[\psi]]v(y,z) = 1
%                                  iff max{[[R(z)]]v[z->a] | a \in |\mathzapf{M}|} = 0, then by our assumptions [[P(x)]]v[x->a] = 1 for all a \in |\mathzapf{M}|
%                                  iff max{[[R(z)]]v[z->a] | a \in |\mathzapf{M}|} = 1, then [[R(z)]]v[z->c] = 1 for some c \in |\mathzapf{M}|, then [[\psi]]v[y ->b,z -> c) = 1 for some some b \in |\mathzapf{M}|

Let: $ \phi = \forall x.P(x) \lor R(x) $, $ \psi = ((\forall y.P(y)) \lor (\exists z.R(z))) $. We show that: $ \phi  \models \psi $:

\begin{subequations}
\begin{align}
[[\phi]] = 1 &\iff [[R(x)]]v[x \longmapsto a] = 1 \text{ or } [[P(x)]]v[x \longmapsto a] = 1 \text{ for all } a \in |\mathzapf{M}|  \nonumber \\
& \iff [[P(x)]]v[x \longmapsto a] = 0 \text{ then } [[R(x)]]v[x \longmapsto a] = 1 \text{ for all } a \in |\mathzapf{M}| \label{eqn:line-4}\\
& \iff [[R(x)]]v[x \longmapsto a] = 0 \text{ then } [[P(x)]]v[x \longmapsto a] = 1 \text{ for all } a \in |\mathzapf{M}| \label{eqn:line-5}\\
  \nonumber
\end{align}
\end{subequations}

The semantic entailment is resolved by the following:
\begin{equation*}
\begin{split} 
  [[\phi]]v(x) &= [[\psi]]v(y,z) = 1 \\
    &\iff [[\forall y.P(y)]]v(y) = 1 \text{ or } [[\exists z.R(z)]]v(z) = 1 
\end{split}
\end{equation*}

Then we go over all the cases:
\paragraph{1)}
if $ [[\forall y.P(y)]]v(y,z) = 1 $, then:
\begin{equation*}
\begin{split} 
  [[\forall y.P(y)]]v(y,z) &= \min\{[[P(y)]]v[y \longmapsto a] | a \in |\mathzapf{M}| \} = 1 \\
  &\iff [[P(y)]]v[y \longmapsto a] = 1 \text{ for all } a \in |\mathzapf{M}| \text{ making } [[\psi]]v[y \longmapsto a] = 1\\
\end{split}
\end{equation*}

\paragraph{2)}
if $ [[\forall y.P(y)]]v(y) = 0 $, then:
\begin{equation*}
\begin{split} 
  [[\forall y.P(y)]]v(y) &= \min\{[[P(y)]]v[y \longmapsto a] | a \in |\mathzapf{M}| \} = 0 \\
  &\iff [[P(y)]]v[y \longmapsto a] = 0 \text{ for some } a \in |\mathzapf{M}| \text { then } [[R(y)]]v[y \longmapsto a] = 1 \\ & \text{ : by assumption \ref{eqn:line-4}, making } [[\psi]]v[y \longmapsto a] = 1
\end{split}
\end{equation*}


\paragraph{3)}
if $ [[\exists z.R(z)]]v(z) = 1 $, then:
\begin{equation*}
\begin{split} 
  [[\exists z.R(z)]]v(z) &= \max\{[[R(z)]]v[z \longmapsto a ] | a \in |\mathzapf{M}| \} = 1 \\
  &\iff [[R(z)]]v[z \longmapsto a] = 1 \text{ for some } a \in |\mathzapf{M}| \text{ making } [[\psi]]v[z \longmapsto a] = 1\\
\end{split}
\end{equation*}

\paragraph{4)}
if $ [[\exists z.R(z)]]v(z) = 0 $, then:
\begin{equation*}
\begin{split} 
  [[\exists z.R(z)]]v(z) &= \max\{[[R(z)]]v[z \longmapsto a ] | a \in |\mathzapf{M}| \} = 0 \\
  &\iff [[R(z)]]v[z \longmapsto a] = 0 \text{ for all } a \in |\mathzapf{M}| \text { then } [[P(y)]]v[y \longmapsto a] = 1 \\ &\text{ : by assumption \ref{eqn:line-5}, making } [[\psi]]v[z \longmapsto a] = 1 \\
\end{split}
\end{equation*}

We've exhausted all cases and all cases lead to  $ [[\phi]]v(x) = [[\psi]]v(y,z) = 1 $, therefore $ \phi  \models \psi $.

\section*{Question 2}
Let: $ \phi = \forall x. \forall y. P(x,y) \rightarrow P(y,x) $, $ |\mathzapf{M}| = \mathbb{N} $, $ \mathzapf{R} = \{P(m,n) = m \dot{=} n \} $, $ \ar(P) = 2 $.

We show that $\phi$ is satisfiable:
\begin{equation*}
\begin{split}
[[\phi]]v(x) &= \min\{[[\phi]]v[x \longmapsto a, y \longmapsto b] | a,b \in |\mathzapf{M}| \} = 1 \\
& = \min\{[[P(x,y)]]v[x \longmapsto a, y \longmapsto b] \implies [[P(y,x)]]v[x \longmapsto a, y \longmapsto b] | a,b \in |\mathzapf{M}| \} = 1 \\
& = \min\{\mathzapf{M}(P)(a,b) \implies \mathzapf{M}(P)(b,a) | a,b \in |\mathzapf{M}| \}  \\
& = \min\{a \dot{=} b \implies b \dot{=} a | a,b \in |\mathzapf{M}| \} \\
& = 1 \text{ : shown by the symmetry property of equality. }
\end{split}
\end{equation*}

This shows that $ \models_M \phi $, therefor $ \phi $ is satisfiable.

\end{document}